\documentclass[11pt]{article}
\usepackage{manual_docs_style}

\begin{document}

\begingroup
\newgeometry{left=0.75in,right=0.75in,top=0.7in,bottom=0.85in}
\pagestyle{plain}
\sffamily
\setlength{\parskip}{0.65em}
\setlist[itemize,1]{leftmargin=1.45em,itemsep=0.15em,topsep=0.25em,parsep=0pt}
\setlist[itemize,2]{leftmargin=2.1em,itemsep=0.08em,topsep=0.08em,parsep=0pt,font=\normalfont}
\setlist[enumerate,1]{leftmargin=1.5em,itemsep=0.2em,topsep=0.25em,parsep=0pt}
\setlist[enumerate,2]{leftmargin=2.1em,itemsep=0.08em,topsep=0.08em,parsep=0pt}
\setlist[enumerate,3]{leftmargin=2.4em,itemsep=0.05em,topsep=0.05em,parsep=0pt}
\makeatletter
\renewcommand\section{\@startsection{section}{1}{\z@}%
  {1.6ex plus .4ex minus .2ex}{0.8ex plus .2ex}%
  {\normalfont\sffamily\Large\bfseries}}
\renewcommand\subsection{\@startsection{subsection}{2}{\z@}%
  {1.2ex plus .3ex minus .2ex}{0.5ex plus .2ex}%
  {\normalfont\sffamily\large\bfseries}}
\makeatother
\newcommand{\mtProjectHeaderImage}{%
  \pdfximage width\textwidth {tex/assets/eth-siplab-thesis-header.png}%
  \noindent\pdfrefximage\pdflastximage\par
}
\mtProjectHeaderImage
\vspace{1.35em}
{\raggedright\Huge\bfseries Deep-learning-based filtering for \name{}\par}
\vspace{0.45em}
{\color{black!15}\hrule height 0.4pt}
\vspace{1.05em}
{\raggedright\large Semester project\par}
\vspace{1.4em}
\phantomsection\label{docstart:learned_validity_filtering_project}
\section*{Research questions}

\begin{itemize}
    \item How much training data do deep-learning time-series models need to obtain high accuracy in \name{}?
    \item How much observational noise can deep-learning time-series models tolerate before classification accuracy degrades?
    \item After calibration, do the model-predicted probabilities accurately reflect empirical correctness?
\end{itemize}

\subsection*{In practical terms}

\begin{itemize}
    \item Can we automate the \name{} filtering process reliably with deep-learning models?
    \item How many samples are rejected by the deep-learning network under noisy conditions?
    \item Can the model’s calibrated confidence score be used as a proxy for sample difficulty?
\end{itemize}

\section*{Introduction}
TraceBench generates controlled root-cause analysis tasks for evaluating whether LLM agents can reason over time-series observations.
It generates these tasks from physical simulators: a system evolves over time, and one parameter may change during the trajectory. 
The agent must infer from the observed time series which parameter changed, or whether no intervention occurred.
For each task to be solvable in principle, each intervention should have a detectable effect on the observed time series. 
Additionally, this effect should be distinguishable from changes to other root-cause parameters.
The current TraceBench pipeline addresses this with a validity-filtering procedure based on probe trajectories, distance-based criteria, and simulator-specific checks.
These rules are interpretable and useful, but they may require manual tuning, introduce human artifacts, and make it difficult to characterize how much noise we can add before interventions are no longer detectable or distinguishable.
This semester project investigates a learning-based alternative. 
Instead of relying only on hand-designed thresholds, the project trains time-series classifiers to score generated samples according to whether the intervention is detectable and distinguishable. 
The training set for the network will be constructed from the output of the \name{} simulation stage, so it may contain ambiguous samples.
We will evaluate both the filtered and unfiltered outputs.
As the test set, we will use a manually cleaned reference set.
The resulting model should be compared against the current filter and, if successful, integrated into the TraceBench benchmark-generation pipeline.

\subsection*{Background}
Learned validity filtering for TraceBench can be cast as a time-series classification (TSC) problem.
CNN-based time-series classifiers can operate directly on raw multivariate trajectories and are strong baselines for TSC \cite{ismailfawaz2019dlreview,wang2017timeseries}. 
InceptionTime is the most established representative of a CNN-based architecture \cite{ismailfawaz2020inceptiontime} and has been widely adopted and benchmarked.
This makes InceptionTime a natural first network for this project.

Recurrent sequence models, including attention-based models and state-space architectures, form a second family of candidates.
Transformers introduced attention-based sequence modeling \cite{vaswani2017attention}, and transformer-based time-series representation learning has been adapted to multivariate time series \cite{zerveas2021transformer}. 
More recent time-series architectures such as TimesNet and PatchTST emphasize multi-period structure and patch-based temporal representations \cite{wu2023timesnet,nie2023patchtst}.
Structured state-space models such as S4 and selective state-space models such as Mamba offer an alternative route to long-context sequence modeling with more favorable scaling than full attention \cite{gu2022s4,gu2023mamba}. 
Released in 2026, MambaSL \cite{jung2026mambasl} is especially relevant because it shows that a minimal single-layer Mamba architecture for time-series classification can be competitive.

Time-series foundation models provide a third, increasingly practical class of candidates.
Forecasting-oriented models such as TimesFM, Chronos, and Lag-Llama pretrain transformer-style models on heterogeneous collections of time series and demonstrate zero-shot or few-shot transfer across datasets, while MOMENT targets more general time-series representation learning and limited-supervision downstream tasks \cite{das2023timesfm,ansari2024chronos,rasul2023lagllama,goswami2024moment}.
For this project, these models can be used as pretrained encoders as their representations may capture generic temporal motifs such as regime shifts, delayed responses, damping behavior, and noise structure.
However, most current time-series foundation models are optimized primarily for forecasting, may make univariate or patched-token assumptions, and may have pretraining biases that are poorly matched to intervention-validity decisions.
Compared with the first two families of algorithms, for this project they represent more of a stretch goal.

Performance should not just be based on accuracy. 
First, we are interested in the confidence, after calibration \cite{guo2017calibration}, so that we can interpret it as probabilities.
Second, we are interested in what changes with more or less noise to characterize task difficulty.


\section*{Task description: Validity filtering}

The goal of this semester project is to design, implement, and evaluate a learning-based validity-filtering component for TraceBench benchmark generation. The method should learn to decide whether generated time-series samples are suitable for inclusion in the benchmark, based on whether the intervention is detectable and distinguishable.

The following tasks and milestones define the project:

\begin{enumerate}[label=\arabic*.]
    \item \textbf{Literature review and codebase familiarization}
    \begin{enumerate}[label=\arabic{enumi}.\arabic*.]
        \item Study TraceBench and its benchmark-generation pipeline
        \begin{itemize}
            \item understand the simulator interface
            \item understand intervention generation
            \item understand labels, metadata, and no-intervention cases
            \item understand the current validity-filtering pipeline
        \end{itemize}
        \item Literature review of time series classification models starting from the background section.
    \end{enumerate}

    \item \textbf{Dataset construction for learned filtering}
    \begin{enumerate}[label=\arabic{enumi}.\arabic*.]
        \item Define data splits: training, validation, and test. 
        Different families must be in different splits.
        \item Standardization and preprocessing (e.g. normalization, downsampling, etc)
    \end{enumerate}

    \item \textbf{Learned validity-filter model}
    \begin{enumerate}[label=\arabic{enumi}.\arabic*.]
        \item Implement a deep-learning classifier that takes a time-series tensor as input and returns class logits, which are converted into calibrated probabilities using a validation set.
        At least InceptionTime, MambaSL, and a transformer-based network should be implemented.
        \item Define augmentation strategies and noise-injection settings.
        \item Log results to Weights \& Biases (wandb) for experiment tracking and lineage.
        Each experiment should be reproducible.
    \end{enumerate}

    \item \textbf{Validity scoring and integration into TraceBench}
    \begin{enumerate}[label=\arabic{enumi}.\arabic*.]
        \item Convert model outputs into filtering decisions
        \begin{itemize}
            \item Generate a curve with the threshold on the x-axis and the number of remaining samples on the y-axis per simulator
            \item Generate a curve with the threshold on the x-axis and accuracy on the y-axis per simulator
            \item Compare these curves across simulators, noise levels, and model families.
        \end{itemize}
        \item Benchmark-generation quality
        \begin{itemize}
            \item number of retained samples per simulator and parameter
            \item diversity of accepted samples
            \item intervention detectability after filtering
            \item distinguishability between possible changed parameters
            \item effect on downstream agent evaluation, if feasible
        \end{itemize}
    \end{enumerate}

    \item \textbf{Analysis and ablations}
    \begin{enumerate}[label=\arabic{enumi}.\arabic*.]
        \item Model ablations:
        \begin{itemize}
            \item Which model performs best?
            \item How does model capacity affect performance?
        \end{itemize}
        \item Data ablations
        \begin{itemize}
            \item Vary the number of training samples
            \item Vary the amount of noise
        \end{itemize}
        \item Failure-case analysis
        \begin{itemize}
            \item inspect accepted invalid samples (relative to the test set)
            \item inspect rejected valid samples (relative to the test set)
        \end{itemize}
        \item Interpretability
        \begin{itemize}
            \item visualize representative trajectories
            \item visualize confidence distributions
            \item optionally use saliency or occlusion analysis to identify influential time windows and variables
        \end{itemize}
    \end{enumerate}

    \item \textbf{Stretch goals}
    \begin{enumerate}[label=\arabic{enumi}.\arabic*.]
        \item Foundation model finetuning
        \item Self-supervised pretraining
        \item Cross-simulator transfer
    \end{enumerate}
\end{enumerate}

\section*{Submission}

\begin{enumerate}[label=\arabic*.]
    \item A written report in the provided two-column template. We recommend adding references and comments on how they relate to the project as they are discovered in .tex.
    \item A clean repository of all code, including dataset generation, training, evaluation, and TraceBench integration.
    The provided GitHub repository needs to be updated every week with notes.
    \item A documented Markdown or .tex file explaining how to replicate the results.
    \item An oral presentation with an interactive Q\&A completes the semester project.
\end{enumerate}

\section*{Mentors}
\begin{itemize}
    \item Tommaso Bendinelli
\end{itemize}

\ifdefined\tracebenchDocsCombinedManualMode
\else
  \bibliographystyle{plain}
  \bibliography{overleaf_paper/references}
\fi

\restoregeometry
\pagestyle{fancy}
\endgroup

\end{document}
